\documentclass[11pt]{article}
\usepackage{amsmath,amssymb,amsthm,hyperref}
\newtheorem{theorem}{Theorem}
\newtheorem{corollary}[theorem]{Corollary}
\newtheorem{proposition}[theorem]{Proposition}
\newcommand{\Qn}{\{0,1\}^n}
\begin{document}
\title{Adaptive Local Covers Require Many Simultaneously Active Denominators}
\author{}
\date{November 23, 2026}
\maketitle

\paragraph{Status and scope.}
This note proves a necessary condition for a specified family of local
construction strategies. It does not prove an improved upper bound on
$W(n)$ or an unrestricted lower bound of order $2^n/n$. In particular,
none of the three requested asymptotic outcomes is asserted here.
The model, head count, and admissibility definitions are unchanged.
The masked scores below are auxiliary mathematical objects, not new
attention heads. Their purpose is to test proposals that first solve
local cells and then discard contributions outside those cells.

\paragraph{What is new relative to the prior audit.}
The programmable-cell obstruction in \texttt{mechanism-audit.tex}
assumes that every labeling of a given cell can be independently
programmed. Here the masks may be selected after seeing the entire
target, and the local pieces may overlap and interact. The parameter
polynomial degree is controlled by the number of pieces that are
simultaneously active at an input. This yields an obstruction even
when no individual cell is universally programmable.

\section{Finite adaptive-mask theorem}

Fix $n,H\geq1$, put $N=2^n$ and $p=1+2H(n+1)$, and let
$\mathcal M$ be a family of $K\geq1$ subsets of $\Qn$.
For an ordered mask list $M_1,\ldots,M_H\in\mathcal M$, define
$I_x=\{h:x\in M_h\}$. Its overlap is
$\max_x|I_x|$. A masked affine-fraction score has the form
\[
 S(x)=c+\sum_{h\in I_x}\frac{A_h(x)}{B_h(x)},
\]
where $A_h,B_h$ are affine and $B_h(x)>0$ throughout $\Qn$.
In particular, all denominators admissible for the actual model are
allowed, but the theorem relaxes their additional slope conditions.
The mask list and all coefficients may depend arbitrarily on the target.

\begin{theorem}[Counting with adaptive masks]\label{thm:count}
Suppose $p\leq N$. The number of strict Boolean sign tables obtainable
from the preceding scores with overlap at most $L$ is at most
\[
 K^H\left(\frac{4e(L+1)N}{p}\right)^p.
\]
Consequently, if all $2^N$ sign tables are obtainable, then
\[
 N\leq H\log_2 K+
 p\log_2\left(\frac{4e(L+1)N}{p}\right).
 \tag{1}\label{eq:entropy}
\]
\end{theorem}
\begin{proof}
First fix a mask list with overlap at most $L$. Use the scalar $c$
and the $2H(n+1)$ affine coefficients as independent real parameters.
For each $x\in\Qn$, multiply $S(x)$ by
$D_x=\prod_{h\in I_x}B_h(x)>0$. The result is
\[
 P_x=c\prod_{h\in I_x}B_h(x)+
 \sum_{h\in I_x}A_h(x)
       \prod_{j\in I_x\setminus\{h\}}B_j(x).
 \tag{2}\label{eq:cleared}
\]
As a polynomial in the parameters, its degree is at most
$|I_x|+1\leq L+1$. The empty active set gives $D_x=1$ and $P_x=c$,
so no endpoint exception is needed. Multiplication by $D_x>0$
preserves both strict positive and strict negative signs.

Warren's strict sign-pattern theorem bounds the number of strict
sign vectors of $N$ polynomials of degree at most $L+1$ in $p\leq N$
variables by $(4e(L+1)N/p)^p$. The coefficient restrictions defining
positive or admissible denominators only remove parameter choices.
There are at most $K^H$ ordered mask lists. There are therefore at most $K^H$ lists satisfying the overlap
condition. Summing the fixed-list bound over those valid lists proves
the claimed bound. Comparing it with $2^N$ and taking binary
logarithms proves \eqref{eq:entropy}.
\end{proof}

The analytic input is H.~E.~Warren,
\emph{Lower Bounds for Approximation by Nonlinear Manifolds},
Transactions of the AMS 133 (1968), 167--178,
\href{https://doi.org/10.2307/1994937}{doi:10.2307/1994937}.
The theorem is used as an established mathematical result, not introduced
as a Lean axiom.

\section{Consequence for every positive exponent improvement}

\begin{corollary}[Exponential overlap is necessary]\label{cor:overlap}
Fix $C,b,\delta>0$. Suppose that for every sufficiently large $n$ all
Boolean sign tables on $\Qn$ have a masked representation satisfying
\[
 1\leq H_n\leq\frac{C2^n}{n^{1+\delta}},
 \qquad K_n\leq2^{bn},
 \qquad \max_x|I_x|\leq L_n.
\]
Then
\[
 \liminf_{n\to\infty}
 \frac{\log_2(L_n+1)}{n^\delta}\geq\frac1{2C}.
 \tag{3}\label{eq:overlap}
\]
Thus a construction with polynomial overlap cannot attain any fixed
positive exponent improvement, even if every center and every local
coefficient is chosen from the complete target table.
\end{corollary}
\begin{proof}
Set $p_n=1+2H_n(n+1)$ and $q_n=2^n/p_n$. Then
\[
 \frac{p_n}{2^n}\leq2^{-n}+
 \frac{2C(n+1)}{n^{1+\delta}}
 =\frac{2C+o(1)}{n^\delta}.
\]
In particular $p_n\leq2^n$ eventually and
$q_n\geq(1/(2C)+o(1))n^\delta$. Dividing
\eqref{eq:entropy} by $p_n$ and rearranging gives
\[
 \log_2(L_n+1)\geq q_n-\log_2q_n-\log_2(4e)
                 -\frac{H_n\log_2K_n}{p_n}.
\]
The final term is at most
$bn/(2(n+1))\leq b/2$. The function
$t-\log_2t$ is increasing for $t>1/\log 2$. Its value at
$(1/(2C)+o(1))n^\delta$ is
$(1/(2C)+o(1))n^\delta$, because $\log n=o(n^\delta)$.
This proves \eqref{eq:overlap}. A polynomial $L_n$ has
$\log_2(L_n+1)=O(\log n)=o(n^\delta)$, giving the stated exclusion.
\end{proof}

\paragraph{Superconstant savings.}
The same proof works with $n^\delta$ replaced by any $g(n)\to\infty$
under $H_n\leq C2^n/(ng(n))$: it gives
$\liminf \log_2(L_n+1)/g(n)\geq1/(2C)$.
Indeed $H_n\geq1$ implies $g(n)\leq C2^n/n$, so
$2^{-n}=o(1/g(n))$, which justifies the same estimate for $p_n/2^n$.
For $g(n)=\log n$, the necessary overlap is polynomial, so the theorem
does not rule out the requested superconstant-saving fallback.

\paragraph{Finite Hamming balls.}
Allow all centers and all radii as masks. This gives
$K_n\leq(n+1)2^n\leq2^{2n}$ for $n\geq1$.
Consequently adaptive ball covers satisfy the dictionary hypothesis.
The exclusion depends on overlap, not on the radius or the number of
points in each ball. If each vertex interacts with only polynomially
many local pieces, selecting a target-dependent cover does not circumvent
Corollary~\ref{cor:overlap}.

\paragraph{Approximate localization.}
Suppose actual scores $F_f$ have strict margin
$\gamma_f=\min_x|F_f(x)|>0$, and a proposed sparsification produces
masked scores $S_f$ of the preceding form with
$\max_x|S_f(x)-F_f(x)|<\gamma_f$.
Then $S_f$ and $F_f$ have identical strict signs at every input: if
$F_f(x)>0$, then $S_f(x)>F_f(x)-\gamma_f\geq0$, and the negative
case is identical with signs reversed. Therefore the same obstruction
applies whenever this sign-preserving sparsification is proved.
Without such a bound the masked theorem says nothing about small tails
of the actual heads; tails can interact through arbitrarily large
cancellations.

\section{Independent cells must grow with the requested saving}

\begin{proposition}[Cell-size requirement at a general rate]
Suppose a cell $E_n\subseteq\Qn$ of size $M_n$ is shattered by sums
of $H_n\geq1$ affine fractions, with positive denominators, and suppose
$H_n\leq CM_n/(ng(n))$, where $C>0$ is fixed and $g(n)\to\infty$.
Then
\[
 \liminf_{n\to\infty}\frac{\log_2 M_n}{g(n)}\geq\frac1{2C}.
\]
In particular, at exponent improvement $g(n)=n^\delta$, an independent
cell must have at least $2^{(1/(2C)+o(1))n^\delta}$ points.
\end{proposition}
\begin{proof}
Set $p_n=1+2H_n(n+1)$. Since $H_n\geq1$, the assumed rate implies
$M_n\geq ng(n)/C$, and hence
\[
 \frac{p_n}{M_n}\leq\frac{C}{ng(n)}+
 \frac{2C(n+1)}{ng(n)}=\frac{2C+o(1)}{g(n)}.
\]
Thus $p_n\leq M_n$ eventually. Clearing denominators and applying
Warren gives
\[
 1\leq\frac{p_n}{M_n}
 \log_2\left(\frac{4e(H_n+1)M_n}{p_n}\right)
 \leq\frac{p_n}{M_n}
 \log_2\left(\frac{4eM_n}{n+1}\right).
\]
The final inequality uses
$(H_n+1)/p_n\leq1/(n+1)$, valid for $H_n\geq1$.
Consequently
$\log_2 M_n\geq M_n/p_n+\log_2(n+1)-\log_2(4e)$.
Divide by $g(n)$ and use the bound on $p_n/M_n$.
\end{proof}

A radius-$r$ Hamming ball has at most $\sum_{j=0}^r\binom nj$
points, at most $(en/r)^r$ when $1\leq r\leq n/2$.
Thus at exponent $\delta$ such an independent ball compiler must have
$r\log_2(en/r)=\Omega(n^\delta)$. In particular fixed radius cannot
suffice. This strengthens the prior constant-head radius-two exclusion
to arbitrary head counts at the proposed improved rate.

\section{Literature transfer and remaining mechanism}

\paragraph{Covering codes.}
Cooper, Ellis, and Kahng,
\href{https://math.iit.edu/~rellis/papers/3asy_short.pdf}
{\emph{Asymmetric Binary Covering Codes}}, JCTA 100 (2002), 232--249,
prove $K^+(n,R)=\Theta_R(2^n/n^R)$ for fixed $R$.
Their theorem transfers directly to a count of downward Hamming-ball
centers. It does not turn a degree-$R$ local numerator into an affine
numerator. The new mask theorem shows that allowing target-dependent
centers and polynomial overlap cannot supply an exponent improvement
through sign-preserving local sparsification. It leaves globally
interacting, nonsparsifiable head systems untouched.

\paragraph{Exact hyperplane covers.}
Aaronson, Groenland, Grzesik, Johnston, and Kielak,
\href{https://arxiv.org/abs/2010.00315}
{\emph{Exact Hyperplane Covers for Subsets of the Hypercube}},
give a combinatorial cover, not an affine-ratio union operation.
Hyperplane membership as a strict sign test does not by itself supply
numerical indicators that can be added. No new transfer was obtained.
The full hyperplane-intersection dictionary is also much larger than
the $2^{O(n)}$ mask dictionaries in Corollary~\ref{cor:overlap}, so that
corollary must not be applied to arbitrary hyperplane masks without
counting their dictionary size.

\paragraph{One concrete missing lemma.}
The new necessary condition suggests the precise type of mechanism
that would avoid this obstruction: an unconditional universal construction
of $O(2^n/(n\log n))$ admissible affine-ratio scores with strict whole-cube
signs, in which the sign information is carried jointly by at least
polynomially many simultaneously significant denominators at some input.
The known two-one-hot majority-ranking compiler is a weaker local tool;
a valid whole-cube gluing lemma with sublogarithmic multiplicative overhead
would turn its logarithmic local saving into the requested fallback.
No such gluing lemma is proved here, and it is not treated as a
conditional solution.

\section{Audit of the separate total-domination branch}

The branch \texttt{ar/1dc3df1f7016} of
\texttt{VerifiedMechanisms/autoresearch-runs} was read at commit
\texttt{045a8fb}. Its exact declaration
\texttt{HeadBound.HStar\_le\_all}, in
\texttt{HeadBound/SubCube.lean}, assumes $1\leq n$ and states
\[
 \texttt{HeadComplexity.HStar}\ n\ f
 \leq 2^n/2^{\texttt{Nat.log}\ 2\ n},
\]
where division is natural-number division. The inspected branch has no
corresponding theorem named against \texttt{worstCaseHeadComplexity}.
The real inequality $W(n)<2^{n+1}/n$ is a direct written consequence;
it is not part of the formal statement of \texttt{HStar\_le\_all}.

The general cover theorem is
\texttt{HeadBound.fracComputable\_of\_sectionCover}, in
\texttt{HeadBound/CoverGeneral.lean}. Its hypotheses provide affine
forms $D_j^0$ whose zero sections cover the cube and, on each section,
affine Lagrange functions that equal one at one chosen vertex and zero
at every other vertex of that section. The proof first establishes full
linear spanning at a degenerate parameter value. The theorem
\texttt{HeadBound.span\_generic}, in \texttt{HeadBound/Generic.lean},
then keeps full rank outside the finitely many roots of a nonzero
polynomial determinant. One chooses an admissible parameter $t>n+1$
and denominators
\[
 D_j(t,x)=D_j^0(x)+t(4n+2|x|).
\]
The parameter $t$ and the denominator list can be chosen independently
of the target. The proof explicitly concludes exact $0/1$ interpolation
before subtracting $1/2$ to obtain strict signs.

Consequently this bridge cannot directly use a sign-only rational
compiler on a radius-two cell. Such a compiler provides neither the
affine Lagrange functions nor the full-rank hypothesis. The determinant
argument preserves linear spanning, not arbitrary sign feasibility
along a path from degenerate to admissible denominators. The prior
mechanism audit already gives a counterexample to that latter inference.
This route remains blocked unless a new admissible sign-gluing theorem
is supplied; the branch's sharper constant does not supply one.

\section{Formal verification status}

The new file
\texttt{formalization/HeadComplexity/Results/ExponentMechanisms.lean}
contains the auxiliary definition \texttt{activeCountingClearedPoly}
and the following kernel-checked statements:
\begin{itemize}
\item \texttt{activeCountingClearedPoly\_totalDegree\_le}: equation
\eqref{eq:cleared} has parameter degree at most $|I_x|+1$.
\item \texttt{eval\_activeCountingClearedPoly}: evaluation equals the
masked score multiplied by its positive active-denominator product.
\item \texttt{strictPolynomialSignPattern\_of\_active\_score}:
every strict masked certificate transfers to the indicated polynomial
sign pattern.
\item \texttt{active\_strictPatterns\_card\_le\_warren}: the fixed-mask
counting consequence, with \texttt{WarrenStrictSignPatternBound} as an
explicit hypothesis.
\end{itemize}
The finite union over mask lists, the asymptotic consequences, and the
cell-size proposition are written proofs only. The project already
isolates Warren's analytic theorem as a hypothesis in
\texttt{CountingLowerBound.lean}; this note does not claim that theorem
has been formalized or remove its hypothesis. No \texttt{sorry}, new
axiom, original definition change, or changed original theorem statement
is introduced.

\end{document}
